\section{Discussion}\label{sec:discussion}
\textbf{To support the claims in \Cref{sec:supervised,sec:downstream}, we run a broad suite of supplementary experiments.} We train all ten OCEAN amplifier and suppressor adapters plus the control on every baseline model, with full per-trait LoRA-scale sweeps on \texttt{TRAIT} and \texttt{MMLU} (\Cref{sec:appendix-cross-model}). On \LlamaThreePointOneSizeEightBInstruct{} we compare alternative DPO recipes (constitution and pair-construction choices; \Cref{sec:appendix-b-dpo-methods}); compress adapter rank (\Cref{sec:downranking}); apply them to base--instruct interpolations (\Cref{sec:instruct_base_interp}); quantify cross-LoRA interactions via pairwise residuals and five-LoRA combinations (\Cref{sec:appendix-residuals-heatmaps,sec:ocean-combinations}); compare against activation capping and system prompting (\Cref{sec:activation_capping,sec:comparison-induction}); and validate PsychAdapter~\citep{vu2026psychadapter} (\Cref{sec:appendix-prefix}).

\textbf{The trait axes carry downstream behavioural signal.} A finding with direct safety relevance is that movement along an OCEAN axis meaningfully changes real-world behaviour. For example, suppressing neuroticism dampens model frustration~\citep{soligo2026gemma}, increasing agreeableness increases sycophancy and compliance with harmful requests (\Cref{sec:downstream}). None of these tasks appear in LoRA training, suggesting that manipulating OCEAN can serve as a lever for tuning a model's behavioural patterns. This is both an opportunity for alignment and also a signal that character training pipelines should be audited for downstream alignment effects. 

\textbf{LLM personas imperfectly reflect human personality structure.} Inter-trait correlations in our LoRAs partly mirror the correlations in human populations (e.g.\ conscientiousness and neuroticism covary in the same direction as in \citet{digman1990personality}) and partly do not.
This is consistent with the broader thesis of the paper: human psychometrics is a useful initialisation for studying LLM behaviour because it comes with validated instruments, but the goal of persona research on LLMs should be to move beyond the human frame and characterise the model's own persona space, and the unsupervised work in \Cref{sec:unsupervised-persona-exploration} is a first step in that direction.

\textbf{Long-term alignment and societal impact.}
More broadly, persona control may be useful because many alignment failures are failures of character: models may be too deferential, too brittle under pressure, too eager to comply, or too willing to preserve a locally rewarded role at the expense of a broader norm. As AI systems become more agentic and are embedded in high-stakes social, scientific, and institutional workflows, their stable behavioural dispositions may shape outcomes across millions of ambiguous interactions
\citep{macaskill2026importance,huang2025valueswilddiscoveringanalyzing}. Our hope is that trait-space methods provide a small step toward a more mature science of model
character and intentions~\citep{marks2026persona,africa2026personas}: first by making behavioural dispositions measurable, then by making them controllable, and eventually by linking them to deeper motivational structure. In this view, persona cartography is a way to audit and shape how models generalise under pressure, helping future systems remain honest and resilient in high-stakes cases where surface behaviour alone is uninformative.

\subsection{Limitations}\label{sec:limitations}

\textbf{Broader models, evaluations, and measurements.} While we train and evaluate adapters across six baseline models spanning a range of sizes and families (\LlamaThreePointOneSizeEightBInstruct, \QwenThreeSizeEightB/\QwenThreeSizeThirtyTwoB, and \GemmaThreeSizeFourBIT/\GemmaThreeSizeTwelveBIT/\GemmaThreeSizeTwentySevenBIT), our full evaluation stack, comprising a psychometric MCQ benchmark, capability benchmarks, and an LLM-judge panel, was applied only to \LlamaThreePointOneSizeEightBInstruct. On the other five models we ran only a subset of this stack: \texttt{TRAIT} and \texttt{MMLU}. This is broader than many prior reports, but it is still a narrow slice of the evaluation space. More fundamentally, \texttt{TRAIT} was designed for humans, and it is not obvious that a human-calibrated questionnaire measures a well-defined construct when applied to an LLM whose persona is not stable across contexts. We may also want to look at personas in multilingual or agentic coding settings. Unlike our trait judges, those for downstream effects are used off-the-shelf without comparable validation.

\textbf{More realistic exploration of traits.} The analysis presented in \Cref{sec:unsupervised-persona-exploration} is exploratory work using highly diverse but nonetheless synthetic rollouts. A study using real deployment trajectories across a range of tasks, and using a range of models, would enable greater confidence in the methodology and possibly more interesting latent dimensions. Validation of trait-induction was performed using the questionnaire which found the factors, rather than independent judges. A separate run of the full discovery pipeline on another model family remains the key untested next step.

\subsection{Related Work}\label{sec:related-work}
\textbf{Prompting persona-like behaviour.}
Persona-like behaviour can be induced through interventions at different points in the model pipeline. The most lightweight way to shape a model persona is through natural-language prompts, which can produce large behavioural changes without modifying model weights~\citep{askell2021general, brown2020language, ouyang2022instructgpt, bai2022constitutional}. Persona prompting and role-play have also been studied as ways of inducing distinct behavioural profiles in LLMs~\citep{deshpande2023toxicity, park2023generative, shanahan2023role}. However, prompting is often brittle: the induced persona may weaken over long rollouts, conflict with later context, or drift as the conversation evolves.

\textbf{Modifying personas at inference time.}
A second class of interventions shapes personas by modifying the model's computation during inference, rather than by changing the prompt or updating the weights~\citep{turner2023activation, zou2023representation, rimsky2024steering, arditi2024refusal, chen2025persona}, and can be more targeted than prompting because they intervene on internal representations associated with particular traits. For persona control, activation capping has been proposed as a way to limit deviations from an assistant-like default~\citep{lu2026assistant}. Nevertheless, steering methods can depend sensitively on layer choice, prompt distribution, rollout length, and the stability of the relevant representation~\citep{tan2024analysing}.

\textbf{Training and post-training to shape model personas.}
A third class of interventions changes the model's weights. Pre-training exposes models to many human and fictional persona-types~\citep{tice2026alignmentpretraining}, while post-training biases this broad simulator toward narrower behavioural defaults, such as helpful, harmless, instruction-following assistant behaviour~\citep{ouyang2022instructgpt, bai2022constitutional, lu2026assistant}. More targeted character-training pipelines fine-tune or adapt models toward particular persona archetypes~\citep{maiya2025opencharacter, li2025big5}.

\textbf{Hybrid approaches which train model modifiers that can be applied at inference time.} A similar approach to ours is PsychAdapter~\citep{vu2026psychadapter}, which trains LoRAs together with per-layer projections that inject a tweakable vector of trait scores into every layer as a dummy token.

\textbf{Task arithmetic is a prior approach for model weight modification.} This method has shown that weight-space directions encode task-specific knowledge and can be composed additively~\citep{ilharco2023editing,huang2024chatvectorsimpleapproach}. We build on these methods by training trait-specific LoRAs and studying whether they provide a more flexible and robust way to construct personas than prompt-only or activation-only control.

